% =========================================================
% Duality and De Morgan
% =========================================================

\subsection{Complement Duality and De Morgan Laws}

\begin{lratopiccard}{Duality and De Morgan --- Structural Signature}
\begin{lrasignaturetable}
\LraSignatureRow{Context}{All complements are relative to a fixed ambient set \(U\).}
\LraSignatureRow{Order reversal}{Complement sends inclusions \(A\subseteq B\) to \(B^c\subseteq A^c\).}
\LraSignatureRow{Binary De Morgan}{Complement exchanges \(A\cup B\) with \(A\cap B\).}
\LraSignatureRow{Indexed De Morgan}{Complement exchanges indexed unions and indexed intersections, with the usual nonempty-index convention for intersections.}
\LraSignatureRow{Duality principle}{Universal identities built from \(\cup,\cap,\varnothing,U\) travel with their duals.}
\LraSignatureRow{Proof move}{Translate to element membership, then exchange \(\neg(\lor)\) with \(\land\neg\) and \(\neg(\exists)\) with \(\forall\neg\).}
\end{lrasignaturetable}
\end{lratopiccard}

\begin{remark*}[Exposition]
Once ordinary and indexed unions and intersections are available, complement
reveals the structural symmetry between the two operations. The order reversal
of complement was proved earlier as an operation-level behavior; here the same
reversal becomes the De Morgan exchange and the general principle of set
duality.
\end{remark*}

\begin{theorembox}{Theorem (De Morgan's Laws)}
\begin{theorem}[De Morgan's Laws]\label{thm:de-morgan}
Let $U$ be a universe and $A,B \subseteq U$. Then
\[
(A \cup B)^c = A^c \cap B^c
\quad\text{and}\quad
(A \cap B)^c = A^c \cup B^c.
\]
\hyperref[prf:de-morgan]{Go to proof.}
\end{theorem}
\end{theorembox}

\begin{lracategoricalcard}{Categorical Rendering --- De Morgan}
\begin{center}
\begin{tikzpicture}[lra categorical diagram]
  \node[lra cat native object] (A) at (0,1.6)
    {\LraCatObject{A}{\{x:x\in A\}}};
  \node[lra cat native object] (B) at (4.8,1.6)
    {\LraCatObject{B}{\{x:x\in B\}}};

  \node[lra cat left object] (join) at (2.4,0.35)
    {\LraCatObject{A\cup B}{\{x:x\in A\lor x\in B\}}};
  \node[lra cat right object] (comp) at (2.4,-1.0)
    {\LraCatObject{(A\cup B)^c}{\{x:x\notin A\land x\notin B\}}};
  \node[lra cat result object] (meet) at (2.4,-2.35)
    {\LraCatObject{A^c\cap B^c}{\{x:x\in A^c\land x\in B^c\}}};

  \draw[lra cat left arrow] (A) -- (join);
  \draw[lra cat left arrow] (B) -- (join);
  \draw[lra cat universal arrow] (join) -- node[lra cat label,right]{complement} (comp);
  \draw[lra cat result arrow] (comp) -- node[lra cat label,right]{same condition} (meet);

  \node[lra cat note] at (2.4,-3.1)
    {complement turns the join condition into the meet of complements};
\end{tikzpicture}

\end{center}
\[
  x\in(A\cup B)^c
  \quad\Longleftrightarrow\quad
  x\in A^c\cap B^c.
\]
\end{lracategoricalcard}

\begin{remark*}[Standard quantified statement]
\[
  \forall A,B\subseteq U,\quad
  (A\cup B)^c=A^c\cap B^c
  \wedge
  (A\cap B)^c=A^c\cup B^c.
\]
\end{remark*}

\begin{remark*}[Predicate reading]
\[
  \operatorname{SetEqual}((A\cup B)^c,A^c\cap B^c)
  \wedge
  \operatorname{SetEqual}((A\cap B)^c,A^c\cup B^c).
\]
\end{remark*}

\begin{remark*}[Interpretation]
Complement exchanges unions with intersections and intersections with unions.
\end{remark*}

\begin{remark*}[Exposition]
De Morgan's laws mirror the propositional logic identities
$\neg(P \lor Q) \equiv \neg P \land \neg Q$ and
$\neg(P \land Q) \equiv \neg P \lor \neg Q$ under the correspondence
$\cap \leftrightarrow \land$, $\cup \leftrightarrow \lor$, and
$A^c \leftrightarrow \neg A$.
\end{remark*}

\LeanFormalizes{thm:de-morgan}{lra-lean}{LRA.VolumeI.Set.Operations.Laws.Duality}{DeMorganUnion}{pending}
\LeanFormalizes{thm:de-morgan}{lra-lean}{LRA.VolumeI.Set.Operations.Laws.Duality}{DeMorganIntersection}{pending}

\begin{dependencies}
\begin{itemize}
  \item \hyperref[def:union]{Union}
  \item \hyperref[def:intersection]{Intersection}
  \item \hyperref[def:complement]{Complement}
\end{itemize}
\end{dependencies}

\begin{theorembox}{Theorem (De Morgan's Laws for Indexed Families)}
\begin{theorem}[De Morgan's Laws for Indexed Families]\label{thm:indexed-de-morgan}
Let $\{A_i\}_{i\in I}$ be an indexed family of subsets of a universe $U$.
Then
\[
U\setminus\bigcup_{i\in I}A_i
=
\bigcap_{i\in I}(U\setminus A_i).
\]
If $I\neq\varnothing$, then
\[
U\setminus\bigcap_{i\in I}A_i
=
\bigcup_{i\in I}(U\setminus A_i).
\]
\hyperref[prf:indexed-de-morgan]{Go to proof.}
\end{theorem}
\end{theorembox}

\begin{remark*}[Standard quantified statement]
\[
  U\setminus\bigcup_{i\in I}A_i
  =
  \bigcap_{i\in I}(U\setminus A_i),
  \qquad
  I\neq\varnothing\to
  U\setminus\bigcap_{i\in I}A_i
  =
  \bigcup_{i\in I}(U\setminus A_i).
\]
\end{remark*}

\begin{remark*}[Interpretation]
Complement exchanges indexed unions with indexed intersections.
\end{remark*}

\LeanFormalizes{thm:indexed-de-morgan}{lra-lean}{LRA.VolumeI.Set.Operations.Laws.Indexed}{ComplementIndexedUnion}{pending}
\LeanFormalizes{thm:indexed-de-morgan}{lra-lean}{LRA.VolumeI.Set.Operations.Laws.Indexed}{ComplementIndexedIntersection}{pending}

\begin{dependencies}
\begin{itemize}
  \item \hyperref[def:indexed-family]{Indexed Family of Sets}
  \item \hyperref[def:indexed-union]{Indexed Union}
  \item \hyperref[def:indexed-intersection]{Indexed Intersection}
  \item \hyperref[def:complement]{Complement}
\end{itemize}
\end{dependencies}

\begin{definitionbox}{Definition (Set Duality)}
\begin{definition}[Set Duality]\label{def:set-duality}
Two set-theoretic expressions over a fixed universe $U$ are \emph{dual} if one
is obtained from the other by simultaneously replacing
$\cup \leftrightarrow \cap$ and $\varnothing \leftrightarrow U$, with
complements unchanged.
\end{definition}
\end{definitionbox}

\begin{remark*}[Standard quantified statement]
\[
  \operatorname{SetDual}(E,E',U)
  \iff
  E'\text{ is obtained from }E\text{ by interchanging }\cup,\cap
  \text{ and }\varnothing,U.
\]
\end{remark*}

\begin{remark*}[Predicate reading]
\[
  \operatorname{SetDual}(E,E',U).
\]
\end{remark*}

\begin{remark*}[Interpretation]
Set duality pairs expressions by exchanging the two basic set operations and
the two extremal sets of a fixed universe.
\end{remark*}

\begin{dependencies}
\begin{itemize}
  \item \hyperref[def:union]{Union}
  \item \hyperref[def:intersection]{Intersection}
  \item \hyperref[def:empty-set]{Empty Set}
  \item \hyperref[def:complement]{Complement}
\end{itemize}
\end{dependencies}

\begin{corollarybox}{Corollary}
\begin{corollary}[\texorpdfstring{\hyperref[prf:set-duality]{Principle of Set Duality}}{Principle of Set Duality}]\label{cor:set-duality}
Any identity involving $\cup$, $\cap$, $\varnothing$, and $U$ that holds for
all subsets of a universe remains valid when each operation and constant is
replaced by its dual.
\hyperref[prf:set-duality]{\textit{Go to proof.}}
\end{corollary}
\end{corollarybox}

\begin{remark*}[Standard quantified statement]
\[
  \forall A_1,\dots,A_n\subseteq U,\quad
  E=E'
  \to
  E^d=(E')^d.
\]
\end{remark*}

\begin{remark*}[Predicate reading]
\[
  \operatorname{SetEqual}(E,E')
  \to
  \operatorname{SetEqual}(E^d,(E')^d).
\]
\end{remark*}

\begin{remark*}[Interpretation]
Any universal identity built from the dual operations has a valid dual identity.
\end{remark*}

\begin{remark*}[Exposition]
To prove a statement involving unions and intersections, it often suffices to
prove one version; the dual statement follows immediately.
\end{remark*}

\begin{remark*}[Exposition]
The two distributive laws
\[
A \cap (B \cup C) = (A \cap B) \cup (A \cap C)
\quad\text{and}\quad
A \cup (B \cap C) = (A \cup B) \cap (A \cup C)
\]
are duals of each other. Proving either one and applying the Principle of Set
Duality yields the other immediately.
\end{remark*}

\begin{dependencies}
\begin{itemize}
  \item \hyperref[def:set-duality]{Set Duality}
  \item \hyperref[ax:extensionality]{Axiom of Extensionality}
\end{itemize}
\end{dependencies}
